\documentclass[11pt,a4paper]{article}
\usepackage[T1]{fontenc}
\usepackage{lmodern,amsmath,amssymb,amsthm,microtype}
\usepackage[margin=27mm]{geometry}
\usepackage[colorlinks=true,linkcolor=blue,urlcolor=blue]{hyperref}
\setlength{\emergencystretch}{2em}
\newtheorem{theorem}{Theorem}[section]
\newtheorem{lemma}[theorem]{Lemma}
\newcommand{\Z}{\mathbb Z}
\newcommand{\Q}{\mathbb Q}
\newcommand{\Pset}{\mathcal P}
\title{A quadratic scaling obstruction for sequential six-unit bases}
\author{}\date{17 November 2026}
\begin{document}\maketitle
\begin{abstract}
We consider two field-valued detectors whose contact weights are $u$ and
$uv$. Four shared comparison tests and four shared reciprocal tests are
distributed independently among their quadratic slots. We prove that
no integer diagonal Laurent scaling makes these factor equations
quadratic using any of $2834$ explicitly specified six-unit bases.
An exact finite enumeration leaves two bases; their exclusion follows
from incompatible weights forced on the reciprocal tests. The result
concerns this factor design and these monomial transformations only.
\end{abstract}

\section{Unit coordinates and allowed transformations}

Let $u,v$ be formal invertible variables. Associate the exponent vector
$(a,b)\in\Z^2$ to the Laurent monomial $u^av^b$. The zero vector
represents the constant $1$. Start with the four unit coordinates
\begin{equation}\label{eq:initial}
u,\quad u^{-1},\quad v,\quad z=u^sv^{-1},
\qquad -2\le s\le2,
\end{equation}
and put
\[
 \Pset_s^{(0)}=\{(0,0),(1,0),(-1,0),(0,1),(s,-1)\}.
\]
The initial guards are $u u^{-1}=1$ and $vz=u^s$. In the second
guard, $u^s$ means $u^2$, $u$, $1$, $u^{-1}$ or $(u^{-1})^2$,
as appropriate, so both guards have degree at most two.

For a finite exponent set $P$ containing zero, define
\[
 \mathcal E(P)=(P+P-P)\setminus P.
\]
Choose $e_1\in\mathcal E(\Pset_s^{(0)})$, then
$e_2\in\mathcal E(\Pset_s^{(0)}\cup\{e_1\})$, and let
\begin{equation}\label{eq:finalP}
 \Pset=\Pset_s^{(0)}\cup\{e_1,e_2\}.
\end{equation}
Thus $\Pset$ describes six unit coordinates and the constant one.
If $e=a+b-c$ is chosen at either stage, its new coordinate $U_e$
has the quadratic definition
\[
                         U_eU_c=U_aU_b.
\]
Every coordinate on the right, and $U_c$, is already available.
The initial guards imply that all initial coordinates are nonzero;
induction gives the same conclusion for both new coordinates.
The definitions therefore determine the new coordinates uniquely.

Every nonunit coordinate $X$ may be replaced by
$u^{w_{X,0}}v^{w_{X,1}}\widetilde X$, with $w_X\in\Z^2$.
Each rational component equation may be divided by an arbitrary
Laurent monomial with integer exponent vector $r$.
The resulting terms must be written literally as ordinary monomials
of total degree at most two in the available coordinates. We allow
neither combinations of factor equations nor substitutions using
other relations among nonunit coordinates.

For a nonunit monomial of degree $d$, the residual unit exponent
must be a sum of at most $2-d$ elements of $\Pset$. In particular,
a quadratic nonunit monomial must have residual exponent zero,
and a constant term has residual exponent in $\Pset+\Pset$.

\section{The two factor systems}

Work in $K=\Q(\rho)$, where $\rho^2=2$, and put
\[
                  q_0=u,\qquad q_1=uv.
\]
There are four shared comparison tests and four shared raw tests.
Use separate field-coordinate blocks
$J_i=J_{i,0}+\rho J_{i,1}$ and
$R_i=R_{i,0}+\rho R_{i,1}$, $1\le i\le4$.
The comparison test is $u^{-6}J_i^2$ and the raw test is $R_i^2$.
Nonzero rational diagonal changes in these component coordinates
are immaterial.

Each detector has four regular factors. Each factor contains a
noncancelled constant $1$ and two quadratic-test slots. For detector
$j$, the allowed contributions are
\begin{equation}\label{eq:slots}
\begin{array}{c|cc}
 &\text{first slot}&\text{second slot}\\\hline
\text{comparison }i&q_j^9u^{-6}J_i^2&q_j^{11}u^{-6}J_i^2\\
\text{raw }i&q_jR_i^2&q_j^3R_i^2.
\end{array}
\end{equation}
The eight tests are assigned bijectively to the eight slots in each
detector. The assignments in the two detectors are independent.
Any remaining terms, including the factor output, are linear in
separate nonunit coordinates and cancel none of the displayed
constant or square monomials.

For the elliptic detector normalization in \cite{Physical}, the
four comparisons are $C_2^4,C_3^4,C_4^4,G^4$. Their stored
squares have the form $u^{-3}J_i$, after fixed nonzero diagonal
normalization; hence their fourth powers have the form used in
\eqref{eq:slots}. The raw tests are reciprocal squares stored
without a fixed Laurent factor. The contact powers in
\eqref{eq:slots} are those of \cite{Factors}.

\begin{lemma}\label{lem:components}
In a quadratic presentation of a factor containing a full square,
the two components of that square have equal Laurent weights.
The real and imaginary factor equations have equal row-division
exponents, even if they were initially allowed independent exponents.
\end{lemma}
\begin{proof}
For a square with unit exponent $e$, the real equation contains
the two nonzero monomials
$u^{e_0}v^{e_1}X_0^2$ and
$2u^{e_0}v^{e_1}X_1^2$. Their quadratic degree forces
\[
                 e+2w_{X_0}=r_{\rm re}=e+2w_{X_1}.
\]
Thus $w_{X_0}=w_{X_1}$. The nonzero imaginary monomial
$2u^{e_0}v^{e_1}X_0X_1$ then gives
$r_{\rm im}=e+w_{X_0}+w_{X_1}=r_{\rm re}$.
The different test blocks prevent cancellation of these terms.
\end{proof}

\section{The finite unit enumeration}

For a comparison occurring with powers $p_0,p_1\in\{9,11\}$,
write $w\in\Z^2$ for its common component weight, and let $r_j$
be the row exponent of the factor containing it in detector $j$.
Lemma~\ref{lem:components} gives
\begin{equation}\label{eq:r}
 r_0=2w+(-6+p_0,0),\qquad
 r_1=2w+(-6+p_1,p_1).
\end{equation}
Since both factors have constant $1$, their constant exponents
$E_j=-r_j$ belong to $\Pset+\Pset$. Consequently
\begin{equation}\label{eq:filter}
\begin{split}
 E_0,E_1&\in\Pset+\Pset,\\
 E_1-E_0&=(p_0-p_1,-p_1),\\
 E_0&\equiv(-6+p_0,0)\pmod2.
\end{split}
\end{equation}

\begin{lemma}\label{lem:enumeration}
Among the bases \eqref{eq:finalP}, exactly two admit a state
\eqref{eq:filter}. They have $s=0$, $e_1=(0,-3)$ and
$e_2=(\varepsilon,-6)$, where $\varepsilon\in\{-1,1\}$.
Their complete state lists are
\begin{equation}\label{eq:state-table}
\begin{array}{c|c|c|c|c}
\varepsilon&(p_0,p_1)&E_0&E_1&w\\\hline
-1&(9,9)&(-1,0)&(-1,-9)&(-1,0)\\
 1&(9,9)&(1,0)&(1,-9)&(-2,0)\\
 1&(11,9)&(-1,0)&(1,-9)&(-2,0).
\end{array}
\end{equation}
\end{lemma}
\begin{proof}
The finite enumeration is given by the following exact procedure.
For each $s=-2,-1,0,1,2$, form $P_0=\Pset_s^{(0)}$.
For every $e_1\in(P_0+P_0-P_0)\setminus P_0$, form
$P_1=P_0\cup\{e_1\}$. For every
$e_2\in(P_1+P_1-P_1)\setminus P_1$, form
$P=P_1\cup\{e_2\}$, retaining each resulting set once for that
$s$. For each $(p_0,p_1)\in\{9,11\}^2$ and every
$E_0\in P+P$, check the two parity conditions and membership of
$E_0+(p_0-p_1,-p_1)$ in $P+P$. The resulting counts are
\[
\begin{array}{r|r|r|r}
s&\text{ordered extension choices}&\text{distinct bases for }s
 &\text{bases passing \eqref{eq:filter}}\\\hline
-2&1140&762&0\\
-1&868&592&0\\
 0&648&458&2\\
 1&868&592&0\\
 2&1140&762&0\\\hline
\text{total}&4664&3166&2.
\end{array}
\]
Removing duplicates between different values of $s$ gives $2834$
distinct final exponent sets. The successful state lists are exactly
\eqref{eq:state-table}. All set operations, membership tests and
parity checks are over integers. The standard-library checker in
Section~\ref{sec:certificate} implements precisely these loops,
reconstructs the complete sets and verifies the displayed lists.

For both surviving bases, the two additional quadratic definitions
can be written
\[
                  vU_{(0,-3)}=(v^{-1})^2,
\]
and, respectively,
\[
 uU_{(-1,-6)}=U_{(0,-3)}^2,
 \qquad u^{-1}U_{(1,-6)}=U_{(0,-3)}^2.
\]
These also verify their sequential admissibility directly.
\end{proof}

\section{Incompatible detector weights}

\begin{theorem}\label{thm:obstruction}
For every unit basis \eqref{eq:finalP}, the two factor systems
\eqref{eq:slots} admit no quadratic presentation by the integer
diagonal Laurent transformations and monomial row divisions specified
above. The conclusion holds for every independent assignment of the
quadratic tests and every choice of the additional linear terms.
\end{theorem}
\begin{proof}
Apply \eqref{eq:filter} to any comparison. By
Lemma~\ref{lem:enumeration}, only the two displayed bases remain.
Apply the same lemma separately to all four comparisons. In either
basis they all have the same component weight, namely $(-1,0)$
for $\varepsilon=-1$ and $(-2,0)$ for $\varepsilon=1$.
Moreover, every comparison has power $9$ in detector $1$.

There are four comparisons and four first slots. Thus the comparisons
occupy all first slots of detector $1$, and the raw tests occupy all
second slots, where their power is $3$. Equation~\eqref{eq:r}
shows that every factor row exponent in this detector is $(1,9)$
for $\varepsilon=-1$, or $(-1,9)$ for $\varepsilon=1$.
If $b$ is the common component weight of a raw test, its square
forces
\[
 2b+(3,3)=(1,9)\quad\hbox{or}\quad
 2b+(3,3)=(-1,9),
\]
respectively. Hence every raw test has second weight component $3$.

In detector $0$, two comparisons cannot share one factor: their
common weight and their slot powers $9,11$ would give row exponents
differing by $(2,0)$. Since there are four comparisons and four
factors, each factor therefore contains exactly one comparison and
one raw test. The comparison term forces the second row-exponent
component to be zero. The raw term, whether its power is $1$ or $3$,
has second exponent component $2\cdot3=6$, since $q_0=u$.
The two quadratic terms cannot have the same row exponent.
This contradiction excludes both surviving bases.
\end{proof}

The theorem uses no reference equation, point-chart equation, special
factor, product-chain equation or particular linear-test allocation.
It imposes no bound on the size of the integer Laurent weights.
Different unit families, detector powers, nonmonomial coordinates,
or operations using other equations are outside its hypotheses.
In particular it gives no lower bound for unrestricted quartic
Diophantine definitions.

\section{Exact certificate}\label{sec:certificate}

The portable specification is \texttt{factor\_model.json};
\texttt{check.py} reconstructs every finite exponent set, every
surviving comparison state, and the vector contradiction in
Theorem~\ref{thm:obstruction}. It requires only the Python standard
library and accepts
\begin{quote}\ttfamily
python -B check.py --output certificate.json
\end{quote}
The certificate records the per-$s$ and global counts, hashes of the
sorted finite sets, both quadratic unit-definition chains, all
surviving states and both incompatible raw-test weights. Its proof
does not depend on an SMT solver.

\begin{thebibliography}{9}\raggedright
\bibitem{Physical}
\emph{Laurent scaling obstructions for two elliptic detectors},
\nolinkurl{negative_results/papers/quartic/physical_detector_scaling_20261117_v67.tex},
with its exact physical-coordinate mapping certificates.
\bibitem{Factors}
\emph{Five local factors with squared reciprocal tests},
\nolinkurl{papers/quartic/revised_detector_factors_20261117_v66.tex}.
\end{thebibliography}
\end{document}
